\section{Solution}

Seen through \citet{seidl2013meta}, the paradox turns on the relative magnitude of two ratios: it arises whenever the growth rate of utility, $u(x_{n+1})/u(x_n)$, exceeds the decay rate of probability, $p(x_n)/p(x_{n+1})$, for infinitely many outcomes.
Any solution must therefore either restrain the growth of utility or restrict the class of lotteries over which the comparison is made.

The first response bounds utility.
\citet{menger1934unbounded} shows that boundedness is necessary as well as sufficient: expected utility is finite for all St.\ Petersburg-type lotteries if and only if there exists a $\bar{u} < \infty$ with $u(x) \leq \bar{u}$ for all $x$, in which case $\mathbb{E}[u(L)] \leq \bar{u} < \infty$ \citep{seidl2013meta}.
\citet{menger1934unbounded} himself rejects boundedness as a description of behaviour, but the restriction turns out to be unavoidable within the axiomatic theories: \citet{fishburn1970utility} shows that the axioms of \citet{savage1954axiom} already imply a bounded utility function.

The second response restricts the set of admissible lotteries rather than the utility function. \citet{arrow1974unbounded} shows that if $u$ is concave and monotone increasing with $u(0)$ finite, and the lottery has a finite mean, then expected utility is finite.
The argument follows from Jensen's inequality, since concavity of $u(\cdot)$ gives
\begin{align*}
    \mathbb{E}[u(L)] \leq u\!\left(\mathbb{E}[L]\right) < \infty,
\end{align*}
where the right-hand side is finite whenever $\mathbb{E}[L] < \infty$ \citep{seidl2013meta}.

Restrictions of this kind have a long pedigree.
\citet{bernoulli1954exposition} reports the proposal to bound the winnings at $2^{24}$ coins, and later authors justified bounded winnings by the finiteness of the bookmaker's wealth or of the time available to play \citep{seidl2013meta}.
A parallel strand introduces finiteness on the probability side instead, treating probabilities below $1/10{,}000$ as zero \citep{seidl2013meta}.